\documentclass[aspectratio=169,10pt]{beamer}
\usepackage[UTF8,fontset=none]{ctex}
\usepackage{booktabs}
\usepackage{array}
\usetheme{nju-showcase}

\title[NJU-Beamer Showcase]{NJU-Beamer Showcase\\南京大学幻灯片模板 · 炫技版}
\subtitle{渐变封面 / 章节导航 / 进度页脚 / 数据组件}
\NJUsetspeaker[汇报人]{NJU-Beamer Team}{南京大学 \quad 某某学院}
\date{2026年11月}
\NJUsetupfootline{NJU-Beamer Showcase}{2026}{NJU-Beamer Team}

\begin{document}
\NJUtitleframe

\section{页面骨架}

\begin{frame}[t]{标题区与页脚}
\begin{columns}[T,onlytextwidth]
  \begin{column}{.56\textwidth}
    标题上方的小字是\NJUmark{章节面包屑}，由 \texttt{\textbackslash section} 自动生成；标题下方是紫到金的渐变分隔线。

    页脚最底部是\NJUmark{全局进度条}，中间的圆点对应全部章节：已讲完的为紫色，当前章节为金色。

    \begin{itemize}
      \item 章节页自带议程：右侧列出全部章节并高亮当前位置
      \item 封面、章节页与尾页共用深紫渐变与丝带曲线
      \item \alert{alert} 为主题紫加粗，\textbf{textbf} 为标准加粗
    \end{itemize}
  \end{column}
  \begin{column}{.40\textwidth}
    \begin{njuquote}{使用方法}
      \texttt{\textbackslash usetheme\{nju-showcase\}}

      其余接口与 \code{nju} 主题完全相同，文稿可以直接切换。
    \end{njuquote}
    \NJUchip{nju} \NJUchip{nju-plain} \NJUchip{nju-report} \NJUchip[NJUCalloutGold]{nju-showcase}
  \end{column}
\end{columns}
\end{frame}

\section{数据与流程}

\begin{frame}[t]{关键数字}
\begin{columns}[T,onlytextwidth]
  \begin{column}{.235\textwidth}\NJUstat{1902}{建校年份}{三江师范学堂}\end{column}
  \begin{column}{.235\textwidth}\NJUstat[NJUCalloutGold]{4}{主题风格}{一套接口}\end{column}
  \begin{column}{.235\textwidth}\NJUstat[NJUCalloutBlue]{16:9}{页面比例}{aspectratio=169}\end{column}
  \begin{column}{.235\textwidth}\NJUstat[NJUCalloutGreen]{100\%}{矢量素材}{校徽、校名、北大楼}\end{column}
\end{columns}

\vspace{4pt}
\texttt{\textbackslash NJUstat[颜色]\{数字\}\{标签\}\{备注\}} 适合放在多栏中并排展示关键数字。

\NJUmeter{0.82}{进度条 \texttt{\textbackslash NJUmeter}}{82\%}
\NJUmeter[NJUCalloutGold]{0.46}{可指定颜色}{46\%}
\end{frame}

\begin{frame}[t]{流程条与提示框}
\NJUsteps{提出问题,建立方法,计算筛选,实验验证,器件集成}

\vspace{6pt}
\begin{columns}[T,onlytextwidth]
  \begin{column}{.32\textwidth}
    \begin{njuquote}{强调引用}
      彩色标题压在一条细线上，正文不带底色和边框。
    \end{njuquote}
    \begin{njuquotegold}{金色提示}
      适合补充说明。
    \end{njuquotegold}
  \end{column}
  \begin{column}{.32\textwidth}
    \begin{njuquoteblue}{蓝色提示}
      适合方法或定义。
    \end{njuquoteblue}
    \begin{njuquotegreen}{绿色提示}
      适合已完成的结果。
    \end{njuquotegreen}
  \end{column}
  \begin{column}{.32\textwidth}
    \begin{njuquotered}{红色提示}
      适合风险与局限。
    \end{njuquotered}
    \begin{block}{Beamer block}
      标准 \code{block} 同样适用。
    \end{block}
  \end{column}
\end{columns}
\end{frame}

\section{图表与代码}

\begin{frame}[t]{图文与表格}
\begin{columns}[T,onlytextwidth]
  \begin{column}{.36\textwidth}
    \begin{njuquote}{图片}
      \begin{Figure}[北大楼剪影：\code{nju-skyline.pdf}]
        \includegraphics[width=.92\linewidth]{\NJUasset{nju-skyline.pdf}}
      \end{Figure}
    \end{njuquote}
  \end{column}
  \begin{column}{.60\textwidth}
    \begin{Table}[四种风格对照]
      \footnotesize
      \begin{tabular}{@{}lccc@{}}
        \toprule
        主题 & 正文背景 & 封面 & 页脚 \\
        \midrule
        \texttt{nju} & 校徽、北大楼、校训 & 斜切紫色 & 页码 \\
        \texttt{nju-plain} & 无 & 斜切紫色 & 三格信息栏 \\
        \texttt{nju-report} & 无 & 居中标题块 & 三格信息栏 \\
        \texttt{nju-showcase} & 校徽水印 & 渐变与丝带 & 进度条 \\
        \bottomrule
      \end{tabular}
    \end{Table}
  \end{column}
\end{columns}
\end{frame}

\begin{frame}[t,fragile]{代码与公式}
\begin{columns}[T,onlytextwidth]
  \begin{column}{.48\textwidth}
    \begin{enumerate}
      \item 编译单份文稿
      \item 严格检查全部文稿
    \end{enumerate}
\begin{lstlisting}[language=bash]
make render INPUT=templates/nju-showcase.tex
make check
\end{lstlisting}
  \end{column}
  \begin{column}{.48\textwidth}
    \begin{njuquoteblue}{行间公式}
      \[
        E_{\mathrm{r}} = \min_{x\in(0,1)}
        \Bigl[E_{\mathrm{hull}}(x) - (1-x)E_A - xE_B\Bigr]
      \]
    \end{njuquoteblue}
  \end{column}
\end{columns}
\end{frame}

\NJUendframe{诚朴雄伟\quad 励学敦行}
\end{document}
